% Standalone landscape table with embedded bibliography. Compile twice with pdflatex.
\documentclass[10pt,a4paper,landscape]{article}
\usepackage[T1]{fontenc}
\usepackage[utf8]{inputenc}
\usepackage[margin=15mm]{geometry}
\usepackage{lmodern,microtype,array,tabularx,booktabs,xurl}
% Extend tabularx to allow page breaks and repeated table headers.
\usepackage{ltablex}
\keepXColumns
\usepackage[hidelinks]{hyperref}
\newcolumntype{Y}[1]{>{\hsize=#1\hsize\linewidth=\hsize\raggedright\arraybackslash}X}
\setlength{\parindent}{0pt}
\setlength{\parskip}{5pt}
\setlength{\tabcolsep}{4pt}
\renewcommand{\arraystretch}{1.18}
\newcommand{\source}[4]{\bibitem{#1} #2. \emph{#3}. \url{#4}.}
\title{Climate Policy Instruments and Scope 3 Decisions}
\author{TRACE3 --- Work Package 4}
\date{Policy review period 2016--2025}
\begin{document}
\maketitle
This inventory consolidates the 2016-2025 timeline, additional 2030-target instruments, selected national carbon-pricing measures in EU member states, and a dedicated Dutch regulations and incentives section. National taxes introduced before 2016 are included where they operated during the review period; this is not an exhaustive inventory of every member state. Repeated milestones for the same instrument are combined. R = regulation or binding legislation (including legally established incentives); I = initiative, strategy or voluntary programme. Affected firms include both directly regulated entities and indirect beneficiaries or supply-chain partners, as indicated. Firm eligibility is summarized, not reproduced as an exhaustive legal test. Dates after 2025 describe schedules associated with the reviewed measures, not a full current-law update. The Scope 3 column contains analytical interpretations of potential decision channels, not claims that each instrument imposes a Scope 3 obligation or has already delivered reductions.

A supplier's operational emissions may enter its customer's Scope 3 purchased-goods inventory. The customer's own purchased electricity is generally Scope 2; outsourced transport can enter Scope 3. Carbon pricing, tax credits or offset purchases alone do not demonstrate a reduction in the buyer's inventory. CSRD/CSDDD scope descriptions follow the historical measures: the December 2025 Omnibus agreement was provisional, with substantive final adoption occurring after this review period.\cite{ghgp,omnibus}

\clearpage
\section*{Global}
\small
\begin{tabularx}{\textwidth}{Y{0.85}Y{1.00}Y{0.85}Y{1.05}Y{1.25}}
\toprule
\textbf{Instrument and type} & \textbf{Purpose and target} & \textbf{Affected firms} & \textbf{Implementation dates} & \textbf{Potential Scope 3 connection} \\
\midrule
\endfirsthead
\toprule
\textbf{Instrument and type} & \textbf{Purpose and target} & \textbf{Affected firms} & \textbf{Implementation dates} & \textbf{Potential Scope 3 connection} \\
\midrule
\endhead
\midrule \multicolumn{5}{r}{Continued on next page} \\ \endfoot \bottomrule \endlastfoot
Paris Agreement [R: treaty]\cite{paris} & Coordinate national mitigation commitments and reporting; pursue temperature goals. & Governments directly; firms indirectly through national implementation. & Entered into force 4 Nov 2016; national policies have separate schedules. & Context for supplier-country climate risk and long-term procurement; no direct corporate Scope 3 mandate. \\
\addlinespace[5pt]
\end{tabularx}
\normalsize

\clearpage
\section*{European Union}
\small
\begin{tabularx}{\textwidth}{Y{0.85}Y{1.00}Y{0.85}Y{1.05}Y{1.25}}
\toprule
\textbf{Instrument and type} & \textbf{Purpose and target} & \textbf{Affected firms} & \textbf{Implementation dates} & \textbf{Potential Scope 3 connection} \\
\midrule
\endfirsthead
\toprule
\textbf{Instrument and type} & \textbf{Purpose and target} & \textbf{Affected firms} & \textbf{Implementation dates} & \textbf{Potential Scope 3 connection} \\
\midrule
\endhead
\midrule \multicolumn{5}{r}{Continued on next page} \\ \endfoot \bottomrule \endlastfoot
Non-financial reporting guidance [I]\cite{nfrdguidance} & Help companies disclose environmental and other non-financial information under the existing NFRD. & Large public-interest entities covered by NFRD; suppliers indirectly through information requests. & Guidance issued Jun 2017; climate supplement Jun 2019; non-binding guidance. & Encourages value-chain information requests and more consistent supplier disclosures. \\
\addlinespace[5pt]
Renewable Energy Directive: RED II / RED III [R]\cite{eu2018,red3} & Raise renewable energy's share: 32\% target in 2018; binding 42.5\% EU target by 2030 in 2023, with 45\% indicative ambition. & Member states directly; energy producers, fuel suppliers and industrial users through implementing rules. & RED II adopted 2018; RED III adopted 2023; national implementation phased; target year 2030. & Cleaner supplier energy can reduce purchased-product footprints; supports supplier renewable-energy engagement. \\
\addlinespace[5pt]
Energy Efficiency Directive [R]\cite{eu2018,eed} & Reduce energy demand; revised target is 11.7\% lower final energy use in 2030 than the 2020 forecast for 2030. & Member states; businesses subject to energy-management/audit requirements through national rules. & Revised 2018 and 2023; phased national implementation; target year 2030. & Supplier efficiency may lower embodied emissions and costs; supports procurement of efficient products. \\
\addlinespace[5pt]
Energy Union Governance Regulation [R]\cite{eu2018} & Coordinate national energy and climate plans and monitor progress. & Member states directly; energy and industrial firms indirectly. & Adopted 2018; plans cover 2021-2030 with updates and reporting. & Provides country-level policy assumptions for supplier location and investment scenarios. \\
\addlinespace[5pt]
European Green Deal [I]\cite{greendeal} & Set the overarching direction toward climate neutrality by 2050. & Firms across sectors indirectly; specific obligations arise from subsequent laws. & Launched Dec 2019; no single company compliance date. & Signals future demand for low-carbon materials, circular products and logistics. \\
\addlinespace[5pt]
Taxonomy Regulation [R]\cite{taxonomy} & Classify environmentally sustainable activities and improve investment transparency. & In-scope reporting companies and financial-market participants; project developers seeking finance. & Entered into force Jul 2020; climate-related application from 2022, with phased disclosures. & May influence financing of supplier upgrades and low-carbon capital goods; no general Scope 3 cap. \\
\addlinespace[5pt]
European Climate Law [R]\cite{climatelaw} & At least 55\% lower EU net GHG emissions by 2030 versus 1990; climate neutrality by 2050. & EU institutions and member states directly; firms through implementing policies. & Entered into force 2021; milestones 2030 and 2050. & Provides a policy scenario for supplier engagement and investment; not a 55\% mandate for each firm. \\
\addlinespace[5pt]
Recovery and Resilience Facility [R: funding]\cite{rrf} & Finance recovery reforms and investment; at least 37\% of national plan allocations support climate objectives. & Eligible project developers, contractors, technology providers and other beneficiaries under national plans. & Established 2021; national projects and milestones implemented through 2026. & Funding may enable supplier retrofits, cleaner transport and low-carbon equipment purchases. \\
\addlinespace[5pt]
CSRD [R]\cite{csrd,stopclock} & Require standardized sustainability reporting and assurance. & Covered large/public-interest firms and other phased categories, including qualifying non-EU groups; suppliers indirectly. & Adopted 2022; first wave reports in 2025 on FY2024; later waves delayed in 2025. & Direct driver of value-chain emissions data collection where material, supplier engagement and target tracking. \\
\addlinespace[5pt]
REPowerEU [I]\cite{repower} & Accelerate energy savings, renewables and diversification away from Russian fossil fuels. & Energy developers, industrial energy users and beneficiaries of national support programmes. & Launched May 2022; implementation through national plans and later measures. & Supports supplier fuel switching and renewables; changes energy-security and sourcing assumptions. \\
\addlinespace[5pt]
EU ETS reform and maritime extension [R]\cite{fit55} & Tighten the carbon cap: 62\% reduction in covered emissions by 2030 versus 2005; extend coverage. & Covered power/industrial installations, aircraft operators and shipping companies; buyers through costs. & Reform adopted 2023; shipping included from 2024, with surrender shares of 40\%, 70\%, 100\% for 2024, 2025, 2026 emissions. & Changes prices of carbon-intensive inputs and freight; encourages low-carbon suppliers and carriers. \\
\addlinespace[5pt]
CBAM [R]\cite{cbam} & Address carbon leakage by aligning the carbon cost of selected imports with EU production. & EU importers of covered cement, iron/steel, aluminium, fertilizers, electricity and hydrogen; non-EU producers supply data. & Reporting began 1 Oct 2023; transitional phase through 2025; financial regime scheduled from 2026. & Prompts supplier-specific embedded-emissions data and sourcing comparisons; CBAM boundaries differ from a full Scope 3 inventory. \\
\addlinespace[5pt]
ESRS, especially E1 [R]\cite{esrs} & Specify sustainability disclosures, including material Scope 1, 2 and 3 emissions, targets and transition information. & Companies reporting under CSRD, subject to materiality and phase-ins. & Adopted 2023; first application to FY2024, reported in 2025; later cohorts follow CSRD schedules. & Direct accounting link: define material value-chain emissions, improve supplier data and disclose progress. \\
\addlinespace[5pt]
CSDDD [R]\cite{csddd,stopclock,omnibus} & Require due diligence on adverse human-rights/environmental impacts; climate-transition planning in the 2024 text. & Large EU firms and qualifying non-EU firms under statutory thresholds; business partners indirectly. & Adopted 2024; no company application in 2016-2025; first phase shifted to 2028 by Apr 2025 law; Dec 2025 reform deal remained provisional. & May support supplier screening and engagement; not a blanket mandate to cut every supplier's GHG emissions. \\
\addlinespace[5pt]
Net-Zero Industry Act [R]\cite{nzia} & Scale clean-technology manufacturing; streamline permitting and apply relevant procurement provisions. & Net-zero technology manufacturers, project developers and public purchasers. & Adopted 2024; implementation phased after entry into force. & Can improve availability of low-carbon capital goods and alter technology/supplier selection. \\
\addlinespace[5pt]
Stop-the-clock directive [R]\cite{stopclock} & Postpone specified CSRD and CSDDD implementation dates. & CSRD second/third waves and firms in the first CSDDD application phase. & Adopted Apr 2025: two-year CSRD delay; one-year CSDDD transposition/first-phase delay. & Changes the timing of mandatory supplier-data and due-diligence programmes; not an emissions target. \\
\addlinespace[5pt]
Clean Industrial Deal and CISAF [I; state-aid framework]\cite{cleandeal} & Support clean energy, industrial decarbonization and clean-technology production. & Eligible industrial firms, clean-tech manufacturers and energy projects receiving national support. & Deal launched Feb 2025; state-aid framework adopted Jun 2025; individual schemes have separate dates. & May improve economics of supplier decarbonization, low-carbon materials and long-term purchasing contracts. \\
\addlinespace[5pt]
Effort Sharing Regulation [R]\cite{esr} & Binding national targets for covered sectors: collective 40\% cut by 2030 versus 2005 after the 2023 revision. & Member states directly; transport, buildings, agriculture, waste and small industry through national measures. & Adopted 2018; strengthened 2023; annual allocations for 2021-2030. & Influences freight, agricultural sourcing and waste-service emissions via national measures. \\
\addlinespace[5pt]
CO2 standards for new cars and vans [R]\cite{cars} & New-car fleet target 55\% lower and new-van target 50\% lower from 2030, relative to 2021 targets. & Manufacturers of new cars/vans; fleet buyers and leasing/logistics firms indirectly. & Revised 2023; these targets apply 2030-2034. & Influences outsourced transport and vehicle choice; affects use-of-sold-products emissions for manufacturers. \\
\addlinespace[5pt]
\end{tabularx}
\normalsize

\clearpage
\section*{EU member states}
\small
\begin{tabularx}{\textwidth}{Y{0.85}Y{1.00}Y{0.85}Y{1.05}Y{1.25}}
\toprule
\textbf{Instrument and type} & \textbf{Purpose and target} & \textbf{Affected firms} & \textbf{Implementation dates} & \textbf{Potential Scope 3 connection} \\
\midrule
\endfirsthead
\toprule
\textbf{Instrument and type} & \textbf{Purpose and target} & \textbf{Affected firms} & \textbf{Implementation dates} & \textbf{Potential Scope 3 connection} \\
\midrule
\endhead
\midrule \multicolumn{5}{r}{Continued on next page} \\ \endfoot \bottomrule \endlastfoot
Sweden: carbon tax [R: tax]\cite{swedentax} & Price fossil-fuel CO2 emissions to encourage efficiency and fuel switching; coverage includes exemptions and interaction with the EU ETS. & Liable fuel suppliers and taxable fossil-fuel users; transport and heating customers through pass-through. & Introduced 1991; continued through the 2016-2025 review period with changes to rates and exemptions. & May change supplier heat/fuel costs and purchased freight emissions; model the tax only for covered uses. \\
\addlinespace[5pt]
France: carbon component of energy excise [R: tax]\cite{francetax} & Add a CO2-based component to fossil-fuel excise duties; reached EUR 44.60/tCO2 in 2018. & Liable suppliers of taxable petroleum products, natural gas and coal; business fuel users indirectly, subject to exemptions. & Introduced 2014; increased during 2016-2018; distinguish the carbon component from total fuel excise. & Affects fuel-intensive suppliers and logistics costs; can improve economics of efficiency and low-carbon transport. \\
\addlinespace[5pt]
Ireland: carbon taxes on fuels [R: tax]\cite{irelandtax,irelandstart} & Price emissions from liable fuels; Finance Act 2020 legislated a path to EUR 100/tCO2 by May 2030. & Liable fuel suppliers and taxable fuel users; transport, heating and industrial consumers indirectly, with reliefs/exemptions. & Natural gas carbon tax began May 2010; 2020 legislation set annual increases through 2030; implementation dates differ by fuel. & Provides an explicit future carbon-cost path for supplier investment, heating and freight decisions. \\
\addlinespace[5pt]
Denmark: industrial green tax reform [R: tax]\cite{denmarktax} & Phase in higher industrial CO2 taxation; headline 2030 rates DKK 750/t outside ETS and DKK 375/t for ETS firms, in 2022 prices, with special treatments. & Covered industrial and energy businesses; treatment varies with ETS participation and activity. & Reform agreed 2022; implementing legislation adopted; phase-in 2025-2030. & Can influence supplier location, material prices and abatement investment; do not apply headline rates uniformly across activities. \\
\addlinespace[5pt]
Germany: national fuel ETS [R: trading, not tax]\cite{germanyprice,germanystart} & Price CO2 from covered heating/transport fuels; fixed certificate prices initially, reaching EUR 55/t in 2025. & Businesses placing covered fuels on the market directly; heating and transport users through price pass-through. & Started 1 Jan 2021; fixed-price phase 2021-2025. & Changes freight and supplier heating costs; separate from carbon taxes and EU ETS when coding policy exposure. \\
\addlinespace[5pt]
\end{tabularx}
\normalsize

\clearpage
\section*{Netherlands}
\small
\begin{tabularx}{\textwidth}{Y{0.85}Y{1.00}Y{0.85}Y{1.05}Y{1.25}}
\toprule
\textbf{Instrument and type} & \textbf{Purpose and target} & \textbf{Affected firms} & \textbf{Implementation dates} & \textbf{Potential Scope 3 connection} \\
\midrule
\endfirsthead
\toprule
\textbf{Instrument and type} & \textbf{Purpose and target} & \textbf{Affected firms} & \textbf{Implementation dates} & \textbf{Potential Scope 3 connection} \\
\midrule
\endhead
\midrule \multicolumn{5}{r}{Continued on next page} \\ \endfoot \bottomrule \endlastfoot
Netherlands: industrial CO2 levy [R: tax]\cite{nltax,nltaxets} & Encourage industrial emissions reductions above the exempt amount; eligible ETS participants deduct the specified ETS price from the levy rate. & Covered industrial installations, including relevant non-ETS installations; customers through input costs. & Effective 1 Jan 2021; annual levy and dispensation-right calculations during the review period; not a tax on every emitted tonne. & Changes supplier abatement economics and potential material prices; account for exempt rights and ETS interaction to avoid double-counting costs. \\
\addlinespace[5pt]
Climate Act (Klimaatwet) [R: framework law]\cite{nlclimate} & Set national climate governance: government pursues 55\% lower GHG emissions by 2030 versus 1990 and climate neutrality by 2050 under the 2023 text. & Government directly; firms through implementing sector rules and incentives, not a uniform company reduction quota. & Original act 2019; strengthened text effective 22 Jul 2023; targets 2030 and 2050. & Provides scenarios for Dutch supplier engagement, sourcing and investment; does not directly impose a firm-level Scope 3 target. \\
\addlinespace[5pt]
SDE++ [I: subsidy programme]\cite{nlsde,nlsdescope} & Support large-scale renewable energy and CO2-reducing technologies by addressing the gap between costs and revenues, subject to scheme rules. & Eligible companies and nonprofit organizations developing supported energy or industrial decarbonization projects. & SDE++ application rounds began autumn 2020; subsequent annual rounds; commissioning and support periods are project-specific. & May enable supplier electrification, low-carbon heat and carbon capture; evaluate actual product-level reductions rather than assuming subsidy receipt reduces Scope 3. \\
\addlinespace[5pt]
Energy Investment Allowance (EIA) [R: tax incentive]\cite{nleia} & Permit extra deduction of qualifying energy-saving investment costs from taxable profit; 40\% deduction in 2025, not a 40\% cash subsidy. & Eligible businesses paying Dutch income or corporate tax and investing in new qualifying Energy List assets. & Long-running annual scheme active during 2016-2025; 2025 percentage and Energy List apply to qualifying 2025 investments. & Improves economics of supplier efficiency and eligible equipment; can lower purchased-product emissions through actual supplier energy savings. \\
\addlinespace[5pt]
MIA and Vamil [R: tax incentives]\cite{nlmia} & MIA allows additional deduction up to 45\% of qualifying investment; Vamil allows flexible depreciation of 75\%, according to asset/list conditions. & Eligible Dutch income/corporate taxpayers investing in assets on the applicable Environmental List. & Annual schemes active during 2016-2025; asset eligibility and benefit depend on investment year; figures describe the 2025 scheme. & Can support supplier circularity and cleaner equipment; check lifecycle effects because environmental eligibility is broader than GHG reduction. \\
\addlinespace[5pt]
Energy-saving and reporting obligations [R]\cite{nlenergy,nlenergydate} & Require applicable energy-saving measures with payback of five years or less, plus reporting or investigation obligations. & Covered business/institutional locations using at least 50,000 kWh electricity or 25,000 m3 natural-gas equivalent annually, subject to scope rules. & Existing duty expanded from Jul 2023; reporting deadline 1 Dec 2023 in that cycle; generally four-year reporting. & Supplier compliance can reduce embodied emissions and reveal investment opportunities; the buyer's own energy savings are principally Scope 1/2. \\
\addlinespace[5pt]
Municipal zero-emission zones for urban freight [R: local access rules]\cite{nlze,nlze2025} & Restrict access by emitting commercial vans and trucks, with transitional arrangements and exemptions. & Carriers, tradespeople, delivery fleets and other van/truck operators entering participating municipal zones. & First zones started Jan 2025; municipal start dates vary; transitional access depends on vehicle class and registration date. & Directly affects carrier selection, last-mile routing and outsourced delivery emissions; a zero-emission tailpipe does not mean zero lifecycle emissions. \\
\addlinespace[5pt]
AanZET [I: purchase subsidy]\cite{nlaanzet,nlaanzetstart} & Support purchase of eligible new zero-emission trucks; funding and support percentages depend on application round and firm size. & Eligible businesses/nonprofits buying qualifying trucks, including eligible leasing companies. & First opened 9 May 2022; subsequent rounds including 2025; procurement and delivery conditions are round-specific. & May lower the cost of contracted zero-emission freight and facilitate buyer-carrier fleet-transition agreements. \\
\addlinespace[5pt]
SEBA [I: purchase subsidy, closed]\cite{nlseba} & Support eligible new zero-emission commercial vans; historical programme rather than an open 2025 incentive. & Eligible entrepreneurs and nonprofit organizations purchasing or financially leasing qualifying vans. & Ran 15 Mar 2021 to 31 Dec 2024; closed to new applications in 2025. & Supported cleaner delivery fleets and may affect later outsourced-transport emissions; do not include it as available funding for new 2025 purchases. \\
\addlinespace[5pt]
\end{tabularx}
\normalsize

\clearpage
\section*{United States}
\small
\begin{tabularx}{\textwidth}{Y{0.85}Y{1.00}Y{0.85}Y{1.05}Y{1.25}}
\toprule
\textbf{Instrument and type} & \textbf{Purpose and target} & \textbf{Affected firms} & \textbf{Implementation dates} & \textbf{Potential Scope 3 connection} \\
\midrule
\endfirsthead
\toprule
\textbf{Instrument and type} & \textbf{Purpose and target} & \textbf{Affected firms} & \textbf{Implementation dates} & \textbf{Potential Scope 3 connection} \\
\midrule
\endhead
\midrule \multicolumn{5}{r}{Continued on next page} \\ \endfoot \bottomrule \endlastfoot
Phase 2 heavy-duty vehicle standards [R]\cite{usphase2} & Reduce fuel consumption and GHG emissions from new trucks and engines. & Heavy-duty vehicle/engine manufacturers; freight operators indirectly. & Adopted 2016; main truck/engine standards phased through model years 2021, 2024 and 2027. & Cleaner carrier fleets may reduce purchased freight emissions; supports carrier-selection criteria. \\
\addlinespace[5pt]
California SB 32 [R]\cite{sb32} & Reduce statewide GHG emissions at least 40\% below 1990 by 2030. & California state agencies directly; firms through implementing measures. & Enacted 2016; target year 2030. & Shapes supplier investment, California sourcing and logistics scenarios; not a uniform firm-level target. \\
\addlinespace[5pt]
California AB 398 / cap-and-trade extension [R]\cite{california} & Extend the existing emissions-trading framework through 2030. & Covered industrial emitters, electricity suppliers/importers and fuel suppliers. & Enacted 2017; extended framework governs the post-2020 period through 2030 as adopted. & Carbon costs can pass into purchased inputs and freight, changing procurement economics. \\
\addlinespace[5pt]
Section 45Q expansion [R: tax incentive]\cite{us45q,ira} & Make eligible carbon capture, utilization and storage investments more attractive. & Eligible capture project owners and partners, including industrial and energy facilities. & Expanded 2018; enhanced again by IRA in 2022; eligibility depends on project/tax rules. & May lower actual supplier product footprints where capture is attributable and verified; credits alone do not establish Scope 3 reductions. \\
\addlinespace[5pt]
AIM Act [R]\cite{aim} & Phase down HFC production/consumption 85\% by 2036 versus statutory baselines. & HFC producers/importers and affected refrigeration, air-conditioning and equipment businesses. & Enacted Dec 2020; allowance phasedown began 2022; sector restrictions phased. & Supports lower-GWP equipment and cold-chain procurement; effects depend on leakage and inventory boundaries. \\
\addlinespace[5pt]
Bipartisan Infrastructure Law [R: funding]\cite{bil} & Fund clean-energy infrastructure and demonstrations, including more than USD 62 billion for DOE. & Eligible utilities, manufacturers, developers, transport projects and infrastructure contractors. & Enacted Nov 2021; awards and deployment occur over multiple programme-specific years. & May enable lower-carbon freight infrastructure and supplier technology adoption. \\
\addlinespace[5pt]
Inflation Reduction Act [R: incentives]\cite{ira,ev2025} & Support clean electricity, manufacturing, hydrogen, vehicles, efficiency and carbon capture. & Eligible project developers, manufacturers, commercial buyers and other tax-credit recipients. & Enacted Aug 2022; provisions phased from 2022/2023 onward; some incentives changed in 2025. & Can lower costs of supplier abatement and low-carbon inputs; sourcing conditions can influence supplier location. \\
\addlinespace[5pt]
Oil and gas methane standards [R]\cite{methane2023,methanedelay} & Reduce methane through equipment standards, monitoring and leak controls. & Covered oil/gas operators; states implement guidelines for existing sources. & Finalized Dec 2023, published Mar 2024; phased compliance; several deadlines extended in 2025. & May reduce upstream emissions of purchased fuels; supports verified low-methane fuel sourcing. \\
\addlinespace[5pt]
Power-plant CO2 standards [R]\cite{power2024,powerrepeal} & Limit emissions from covered existing coal and new gas generating units. & Covered power generators; states implement existing-source plans. & Finalized Apr 2024; unit-specific requirements extend into the 2030s; repeal proposed in 2025. & Supplier grid decarbonization may lower embodied product emissions; a buyer's own purchased electricity is generally Scope 2. \\
\addlinespace[5pt]
2025 clean-vehicle credit termination [R]\cite{ev2025} & End specified federal clean-vehicle tax incentives under the 2025 tax law. & New/used vehicle buyers, commercial fleet purchasers, manufacturers and dealers. & Act enacted Jul 2025; affected credits unavailable for vehicles acquired after 30 Sep 2025. & May weaken fleet-electrification economics for logistics suppliers; affects the timing/cost of freight decarbonization. \\
\addlinespace[5pt]
Waste Emissions Charge and rollback [R]\cite{wec} & Originally price excess oil/gas methane; later remove the implementing charge rule. & Covered oil/gas facilities meeting reporting and excess-emissions conditions. & Authorized 2022; implementing rule finalized 2024, disapproved and removed in 2025; do not code as continuously operational. & Changes the expected incentive to reduce upstream fuel emissions; distinct from the methane standards above. \\
\addlinespace[5pt]
\end{tabularx}
\normalsize

\clearpage
\section*{Asia}
\small
\begin{tabularx}{\textwidth}{Y{0.85}Y{1.00}Y{0.85}Y{1.05}Y{1.25}}
\toprule
\textbf{Instrument and type} & \textbf{Purpose and target} & \textbf{Affected firms} & \textbf{Implementation dates} & \textbf{Potential Scope 3 connection} \\
\midrule
\endfirsthead
\toprule
\textbf{Instrument and type} & \textbf{Purpose and target} & \textbf{Affected firms} & \textbf{Implementation dates} & \textbf{Potential Scope 3 connection} \\
\midrule
\endhead
\midrule \multicolumn{5}{r}{Continued on next page} \\ \endfoot \bottomrule \endlastfoot
Japan: climate mitigation tax [R]\cite{japantax} & Price fossil-fuel CO2 emissions; final scheduled rate JPY 289/tCO2. & Fossil-fuel importers/extractors directly; energy users through price pass-through. & Tax began 2012; final planned rate reached Apr 2016. & Changes supplier energy costs and incentives for efficiency/fuel switching; affects embodied input emissions. \\
\addlinespace[5pt]
China: national ETS [R; earlier development initiative]\cite{china2017,china2025} & Improve emissions efficiency through allowance compliance and trading. & Power generators initially; steel, cement and aluminium smelting added in 2025. & Development plan 2017; national trading Jul 2021; sector expansion 2025. & Encourages emissions data and abatement in material suppliers; may affect input prices and sourcing. \\
\addlinespace[5pt]
Singapore: Carbon Pricing Act / tax [R]\cite{sgact,sgtax,sgscope} & Price emissions: SGD 5/tCO2e initially, SGD 25 in 2024-2025; announced SGD 50-80 ambition by 2030. & Facilities emitting at least 25,000 tCO2e annually; customers indirectly. & Act 2018; tax Jan 2019; increase Jan 2024. The 2030 range was an announced trajectory, not a fixed enacted 2030 rate. & Affects carbon-intensive suppliers' costs and abatement; can inform supplier investment and contract decisions. \\
\addlinespace[5pt]
China: carbon peak and neutrality commitment [I]\cite{chinapledge} & Peak CO2 before 2030; carbon neutrality before 2060. & Government bodies and firms across sectors through subsequent implementation. & Announced Sep 2020; target years 2030 and 2060. & Long-term supplier transition and location signal; not a direct corporate Scope 3 obligation. \\
\addlinespace[5pt]
Japan: Green Growth Strategy [I]\cite{japangrowth} & Coordinate industrial and technology pathways toward carbon neutrality by 2050. & Priority-sector manufacturers, energy firms, transport and technology developers. & Launched Dec 2020; updated 2021; programmes have separate schedules. & Supports supplier technology roadmaps and procurement of lower-carbon materials/equipment. \\
\addlinespace[5pt]
Japan: Green Innovation Fund [I: funding]\cite{japanfund} & Support long-term R\&D, demonstration and deployment of decarbonization technologies. & Selected companies and consortia receiving project support. & Established under FY2020 supplementary budget; projects launched from 2021, with multi-year support. & Can support supplier pilots and future low-carbon inputs; funding is not proof of realized reductions. \\
\addlinespace[5pt]
South Korea: Carbon Neutrality Act [R]\cite{korea,koreatarget} & Establish climate governance; 40\% national GHG reduction by 2030 versus 2018 under the implementing framework. & Government bodies directly; firms affected through sector policies and support measures. & Enacted 2021; act/decree effective Mar 2022; target year 2030. & Guides supplier transition expectations and engagement in Korean manufacturing supply chains. \\
\addlinespace[5pt]
India: Energy Conservation Amendment Act [R]\cite{indiaccts} & Enable carbon-credit trading and related energy-transition requirements. & Designated energy-intensive businesses and other entities covered by implementing rules. & Enacted Dec 2022; implementation through later schemes and rules. & Creates the legal basis for supplier carbon-cost and data obligations; no immediate blanket Scope 3 requirement. \\
\addlinespace[5pt]
India: Carbon Credit Trading Scheme [R]\cite{indiaccts} & Create compliance and offset mechanisms; compliance uses GHG-intensity targets. & Designated obligated industrial entities; eligible voluntary offset project developers. & Notified Jun 2023; amended Dec 2023; implementation and targets phased by sector. & May encourage supplier intensity reductions and measurement; offset certificates do not automatically reduce a buyer's Scope 3 inventory. \\
\addlinespace[5pt]
India: National Green Hydrogen Mission [I: incentives]\cite{indiah2} & Support hydrogen production and electrolyser manufacturing; initial INR 19,744 crore outlay. & Selected hydrogen producers, electrolyser manufacturers and participating industrial users. & Approved Jan 2023; support through programme-specific tenders and awards. & May enable lower-carbon fertilizer, steel and chemical inputs; compare lifecycle emissions and delivered cost. \\
\addlinespace[5pt]
Japan: GX Promotion Act [R]\cite{gxact,gxdates} & Finance industrial transition and introduce phased carbon pricing. & Supported industrial projects; large emitters and fossil-fuel suppliers under later mechanisms. & Enacted 2023; ETS full operation scheduled FY2026; fossil-fuel surcharge FY2028. & Changes supplier investment and future carbon-cost expectations; may support long-term low-carbon purchasing contracts. \\
\addlinespace[5pt]
Japan: GX League trial ETS [I: voluntary]\cite{gxleague} & Pilot company target-setting and emissions trading before the later mandatory phase. & Voluntarily participating companies. & Trial ETS began Apr 2023; trial phase FY2023-FY2025. & Provides supplier engagement and emissions-performance signals; participation alone does not establish lower product emissions. \\
\addlinespace[5pt]
Taiwan: carbon fee [R]\cite{taiwan,twscope} & Price covered emissions; standard NTD 300/tCO2e, with preferential rates for approved reduction plans. & Covered large manufacturers and electricity generators. & Rates announced Oct 2024; applies to 2025 emissions, first payment scheduled May 2026. & Encourages supplier reduction plans and may alter input prices; useful for material/electronics supplier discussions. \\
\addlinespace[5pt]
\end{tabularx}
\normalsize
\clearpage
\begin{thebibliography}{99}
\small
References reused from the timeline were accessed on 24 September 2026; supplementary references were checked on 25 September 2026.

\source{ghgp}{GHG Protocol}{Calculation Tools FAQ}{https://ghgprotocol.org/calculation-tools-faq}
\source{omnibus}{Council of the European Union (2025)}{Council and Parliament strike a deal to simplify sustainability reporting and due diligence requirements and boost EU competitiveness}{https://www.consilium.europa.eu/en/press/press-releases/2025/12/09/council-and-parliament-strike-a-deal-to-simplify-sustainability-reporting-and-due-diligence-requirements-and-boost-eu-competitiveness/}
\source{paris}{UNFCCC}{The Paris Agreement}{https://unfccc.int/process-and-meetings/the-paris-agreement}
\source{nfrdguidance}{European Commission (2017; updated 2019)}{Commission guidelines on non-financial reporting}{https://finance.ec.europa.eu/publications/commission-guidelines-non-financial-reporting_en}
\source{eu2018}{Council of the European Union (2018)}{Energy efficiency, renewables, governance of the Energy Union: Council signs off on 3 major clean energy files}{https://www.consilium.europa.eu/en/press/press-releases/2018/12/04/energy-efficiency-renewables-governance-of-the-energy-union-council-signs-off-on-3-major-clean-energy-files/pdf/}
\source{red3}{Council of the European Union (2023)}{Renewable energy: Council adopts new rules}{https://www.consilium.europa.eu/en/press/press-releases/2023/10/09/renewable-energy-council-adopts-new-rules/}
\source{eed}{Council of the European Union (2023)}{Council adopts energy efficiency directive}{https://www.consilium.europa.eu/en/press/press-releases/2023/07/25/council-adopts-energy-efficiency-directive/}
\source{greendeal}{European Commission (2019)}{Communication on the European Green Deal}{https://commission.europa.eu/publications/communication-european-green-deal_en}
\source{taxonomy}{European Commission}{EU taxonomy for sustainable activities}{https://finance.ec.europa.eu/sustainable-finance/tools-and-standards/eu-taxonomy-sustainable-activities_en}
\source{climatelaw}{European Commission}{European Climate Law}{https://climate.ec.europa.eu/areas-action/european-climate-law_en}
\source{rrf}{European Commission}{Finance and the Green Deal}{https://commission.europa.eu/strategy-and-policy/priorities-2019-2024/european-green-deal/finance-and-green-deal_en}
\source{csrd}{Council of the European Union (2022)}{Council gives final green light to corporate sustainability reporting directive}{https://www.consilium.europa.eu/en/press/press-releases/2022/11/28/council-gives-final-green-light-to-corporate-sustainability-reporting-directive/}
\source{stopclock}{Council of the European Union (2025)}{Simplification: Council gives final green light on the stop-the-clock mechanism}{https://www.consilium.europa.eu/en/press/press-releases/2025/04/14/simplification-council-gives-final-green-light-on-the-stop-the-clock-mechanism-to-boost-eu-competitiveness-and-provide-legal-certainty-to-businesses/}
\source{repower}{European Commission}{REPowerEU}{https://commission.europa.eu/topics/energy/repowereu_en}
\source{fit55}{Council of the European Union (2023)}{Fit for 55: Council adopts key pieces of legislation delivering on 2030 climate targets}{https://www.consilium.europa.eu/en/press/press-releases/2023/04/25/fit-for-55-council-adopts-key-pieces-of-legislation-delivering-on-2030-climate-targets/}
\source{cbam}{European Commission}{Carbon Border Adjustment Mechanism}{https://taxation-customs.ec.europa.eu/carbon-border-adjustment-mechanism_en}
\source{esrs}{European Commission (2023)}{Commission Delegated Regulation (EU) 2023/2772 on sustainability reporting standards}{https://eur-lex.europa.eu/legal-content/en/TXT/?uri=CELEX:32023R2772}
\source{csddd}{Council of the European Union (2024)}{Corporate sustainability due diligence: Council gives its final approval}{https://www.consilium.europa.eu/en/press/press-releases/2024/05/24/corporate-sustainability-due-diligence-council-gives-its-final-approval/}
\source{nzia}{Council of the European Union (2024)}{Industrial policy: Council gives final approval to the net-zero industry act}{https://www.consilium.europa.eu/en/press/press-releases/2024/05/27/industrial-policy-council-gives-final-approval-to-the-net-zero-industry-act/}
\source{cleandeal}{European Commission (2025)}{Clean Industrial Deal State Aid Framework (CISAF)}{https://competition-policy.ec.europa.eu/about/contribution-clean-just-and-competitive-transition/clean-industrial-deal-state-aid-framework-cisaf_en}
\source{esr}{Council of the European Union (2023)}{Fit for 55 package: Council adopts regulations on effort sharing and land use and forestry sector}{https://www.consilium.europa.eu/en/press/press-releases/2023/03/28/fit-for-55-package-council-adopts-regulations-on-effort-sharing-and-land-use-and-forestry-sector/}
\source{cars}{Council of the European Union (2023)}{Fit for 55: Council adopts regulation on CO2 emissions for new cars and vans}{https://www.consilium.europa.eu/en/press/press-releases/2023/03/28/fit-for-55-council-adopts-regulation-on-co2-emissions-for-new-cars-and-vans/}
\source{usphase2}{US Environmental Protection Agency (2016)}{Final Rule for Phase 2 Greenhouse Gas Emissions Standards and Fuel Efficiency Standards for Medium- and Heavy-Duty Engines and Vehicles}{https://www.epa.gov/regulations-emissions-vehicles-and-engines/final-rule-phase-2-greenhouse-gas-emissions-standards}
\source{sb32}{California Air Resources Board}{GHG 1990 Emissions Level and 2020 Limit}{https://ww2.arb.ca.gov/ghg-2020-limit}
\source{california}{California Air Resources Board (2018)}{CARB provides further flexibility and cost-containment for companies in cap-and-trade program}{https://ww2.arb.ca.gov/news/carb-provides-further-flexibility-and-cost-containment-companies-cap-and-trade-program}
\source{us45q}{US Department of Energy (2020)}{DOE Loan Guarantees and IRS Guidance on 45Q Tax Credit Could Benefit Carbon Capture Projects}{https://www.energy.gov/edf/articles/doe-loan-guarantees-irs-guidance-45q-tax-credit-could-benefit-carbon-capture-projects}
\source{ira}{US Internal Revenue Service}{Credits and deductions under the Inflation Reduction Act of 2022}{https://www.irs.gov/credits-and-deductions-under-the-inflation-reduction-act-of-2022}
\source{aim}{US Environmental Protection Agency}{Frequent Questions on the Phasedown of Hydrofluorocarbons}{https://www.epa.gov/hfcs/frequent-questions-phasedown-hydrofluorocarbons}
\source{bil}{US Department of Energy (2021)}{President Biden Signs Historic Bipartisan Infrastructure Deal to Usher in the Clean Energy Future}{https://www.energy.gov/policy/articles/president-biden-signs-historic-bipartisan-infrastructure-deal-usher-clean-energy}
\source{ev2025}{US Internal Revenue Service}{Clean vehicle tax credits}{https://www.irs.gov/clean-vehicle-tax-credits}
\source{methane2023}{US Environmental Protection Agency (2023)}{EPA's Final Rule to Reduce Methane and Other Harmful Pollution from Oil and Natural Gas Operations and Related Actions}{https://www.epa.gov/controlling-air-pollution-oil-and-natural-gas-operations/epas-final-rule-reduce-methane-and-other}
\source{methanedelay}{US Environmental Protection Agency (2025)}{2025 Interim Final Rule to Extend Compliance Deadlines}{https://www.epa.gov/controlling-air-pollution-oil-and-natural-gas-operations/2025-interim-final-rule-extend-compliance}
\source{power2024}{US Environmental Protection Agency (2024)}{Final Carbon Pollution Standards: Overview Fact Sheet}{https://www.epa.gov/system/files/documents/2024-04/cps-111-fact-sheet-overview.pdf}
\source{powerrepeal}{US Environmental Protection Agency (2025)}{EPA Proposes Repeal of Regulations for Power Plants}{https://www.epa.gov/newsreleases/epa-proposes-repeal-biden-harris-epa-regulations-power-plants-which-if-finalized-would}
\source{wec}{US Environmental Protection Agency}{Methane Emissions Reduction Program and GHGRP Subpart W}{https://www.epa.gov/inflation-reduction-act/methane-emissions-reduction-program-and-ghgrp-subpart-w-petroleum-and}
\source{japantax}{Ministry of the Environment, Japan}{Tax for Climate Change Mitigation}{https://www.env.go.jp/content/900453366.pdf}
\source{china2017}{National Development and Reform Commission, China (2017)}{Notice issuing the national carbon emissions trading market construction plan for the power-generation sector [Chinese]}{https://zfxxgk.ndrc.gov.cn/web/iteminfo.jsp?id=2944}
\source{china2025}{Ministry of Ecology and Environment, China (2025)}{Questions and answers on the work plan for national carbon-market coverage of steel, cement, and aluminium smelting [Chinese]}{https://www.mee.gov.cn/ywdt/zbft/202503/t20250326_1104767.shtml}
\source{sgact}{Singapore Statutes Online}{Carbon Pricing Act 2018}{https://sso.agc.gov.sg/Act/CPA2018}
\source{sgtax}{Ministry of Sustainability and the Environment, Singapore}{Carbon Pricing Act}{https://www.mse.gov.sg/policies/climate-change/carbon-pricing-act/}
\source{sgscope}{Ministry of Sustainability and the Environment, Singapore}{Climate Game Changer}{https://www.mse.gov.sg/policies/climate-change/climate-game-changer/}
\source{chinapledge}{State Council of the People's Republic of China (2021)}{China's new five-year blueprint paves way for 2060 carbon-neutrality}{https://english.www.gov.cn/news/topnews/202103/09/content_WS6046cf92c6d0719374afa6a5.html}
\source{japangrowth}{Ministry of Economy, Trade and Industry, Japan}{Green Growth Strategy Through Achieving Carbon Neutrality in 2050}{https://www.meti.go.jp/english/policy/energy_environment/global_warming/ggs2050/index.html}
\source{japanfund}{Ministry of Economy, Trade and Industry, Japan (2021)}{Green Growth Strategy Through Achieving Carbon Neutrality in 2050: Overview}{https://www.meti.go.jp/english/policy/energy_environment/global_warming/ggs2050/pdf/ggs_overview_all.pdf}
\source{korea}{Ministry of Environment, Republic of Korea (2022)}{Press release on the entry into force of the Carbon Neutrality Act and enforcement decree, 25 March}{https://www.me.go.kr/eng/web/board/read.do?boardId=1516150&boardMasterId=522&menuId=461}
\source{koreatarget}{Ministry of Environment, Republic of Korea (2022)}{Legislating the carbon-neutrality vision and greenhouse-gas reduction commitment [Korean]}{https://me.go.kr/home/web/board/read.do?boardId=1515170&boardMasterId=1}
\source{indiaccts}{Press Information Bureau, Government of India (2024)}{Government statement on the Indian Carbon Market and Carbon Credit Trading Scheme, 9 December}{https://www.pib.gov.in/PressReleasePage.aspx?PRID=2082528&lang=1&reg=3}
\source{indiah2}{Press Information Bureau, Government of India (2023)}{Cabinet approves National Green Hydrogen Mission}{https://www.pib.gov.in/PressReleaseIframePage.aspx?PRID=1888547&lang=2&reg=48}
\source{gxact}{Ministry of Economy, Trade and Industry, Japan}{GX Promotion Strategy}{https://www.meti.go.jp/policy/energy_environment/global_warming/transition/pathways_to_green_transformation_eng.pdf}
\source{gxdates}{Agency for Natural Resources and Energy, Japan (2025)}{GX Policy: Achieving Decarbonization and Economic Growth Together}{https://www.enecho.meti.go.jp/en/category/special/article/detail_214.html}
\source{gxleague}{Agency for Natural Resources and Energy, Japan (2023)}{Japan's energy policy toward achieving GX (Part 2)}{https://www.enecho.meti.go.jp/en/category/special/article/detail_179.html}
\source{taiwan}{Ministry of Environment, Taiwan (2024)}{Taiwan's Ministry of Environment Announces Fee-Charging Rates of Carbon Fees}{https://www.moenv.gov.tw/en/news/press-releases/31552.html}
\source{twscope}{Ministry of Environment, Taiwan (2023)}{Power Generators and Large Manufacturers Targeted for Carbon Fee Collection for Now}{https://www.moenv.gov.tw/en/news/press-releases/12023.html}
\source{swedentax}{Government Offices of Sweden}{Sweden's carbon tax}{https://www.government.se/government-policy/taxes-and-tariffs/swedens-carbon-tax/}
\source{francetax}{Ministry for Ecological Transition, France}{Fiscalite carbone}{https://www.ecologie.gouv.fr/politiques-publiques/fiscalite-carbone}
\source{irelandtax}{Houses of the Oireachtas (2025)}{Tax Code: Parliamentary Question 346, 22 January 2025}{https://www.oireachtas.ie/en/debates/question/2025-01-22/346/}
\source{irelandstart}{Revenue Commissioners, Ireland}{Natural Gas Carbon Tax}{https://www.revenue.ie/en/companies-and-charities/excise-and-licences/energy-taxes/natural-gas-carbon-tax/index.aspx}
\source{nltax}{Government of the Netherlands}{CO2-heffing voor industrie}{https://www.rijksoverheid.nl/onderwerpen/klimaatverandering/co2-heffing-voor-industrie}
\source{nltaxets}{Dutch Emissions Authority (2022)}{NEa allocates industrial carbon-levy dispensation rights}{https://www.emissieautoriteit.nl/actueel/nieuws/2022/04/28/nea-stort-dispensatierechten-voor-co2-heffing-industrie}
\source{denmarktax}{Ministry of Climate, Energy and Utilities, Denmark (2025)}{Klimaprogram 2025}{https://www.kefm.dk/Media/638948087473597195/Klimaprogram\%202025.pdf}
\source{germanyprice}{German Environment Agency}{Der Nationale Emissionshandel}{https://www.umweltbundesamt.de/der-nationale-emissionshandel}
\source{germanystart}{German Environment Agency (2020)}{CO2 pricing for emissions in heating and transport sectors to start in new year in Germany}{https://www.umweltbundesamt.de/en/press/pressinformation/co2-pricing-for-emissions-in-heating-transport}
\source{nlclimate}{Government of the Netherlands (2023)}{Klimaatwet: version effective 22 July 2023}{https://wetten.overheid.nl/BWBR0042394/2023-07-22/}
\source{nlsde}{Netherlands Enterprise Agency}{SDE++: Facts and figures}{https://www.rvo.nl/subsidies-financiering/sde/aanvragen/feiten-en-cijfers}
\source{nlsdescope}{Netherlands Enterprise Agency}{Stimulation of Sustainable Energy Production and Climate Transition}{https://english.rvo.nl/subsidies-financiering/sde}
\source{nleia}{Dutch Tax Administration (2025)}{Energie-investeringsaftrek 2025}{https://www.belastingdienst.nl/wps/wcm/connect/bldcontentnl/belastingdienst/zakelijk/winst/inkomstenbelasting/verandering_inkomstenbelasting_vorige_jaren/veranderingen-inkomstenbelasting-2025/investeringsaftrek-2025/energie-investeringsaftrek-2025}
\source{nlmia}{Netherlands Enterprise Agency (2025)}{MIA and Vamil tax scheme}{https://english.rvo.nl/subsidies-financing/mia-vamil/mia-and-vamil-tax-scheme}
\source{nlenergy}{Netherlands Enterprise Agency (2023)}{Information on the energy-saving obligation and reporting obligations}{https://www.rvo.nl/onderwerpen/energiebesparingsplicht-2023/voorlichtingsmateriaal}
\source{nlenergydate}{Netherlands Enterprise Agency}{About the energy-saving obligation}{https://www.rvo.nl/onderwerpen/energiebesparingsplicht/over-de-energiebesparingplicht}
\source{nlze}{Netherlands Enterprise Agency / Business.gov.nl}{Zero-emission zones in the Netherlands}{https://business.gov.nl/sustainable-business/sustainable-business-operations/zero-emission-zones-in-the-netherlands/}
\source{nlze2025}{Government of the Netherlands (2024)}{Road Traffic Signs and Regulations in the Netherlands: Article 86e}{https://www.government.nl/site/binaries/site-content/collections/documents/2024/02/09/road-traffic-signs-and-regulations-in-the-netherlands/Road\%2BTraffic\%2BSigns\%2Band\%2BRegulations\%2Bin\%2Bthe\%2BNetherlands.pdf}
\source{nlaanzet}{Netherlands Enterprise Agency}{Aanschafsubsidie Zero-Emissie Trucks (AanZET)}{https://www.rvo.nl/subsidies-financiering/aanzet}
\source{nlaanzetstart}{Netherlands Enterprise Agency (2023)}{Economic significance of the Dutch electric transport sector 2020-2022}{https://www.rvo.nl/sites/default/files/2023-11/RVO\%20-\%20Rapport\%20Economische\%20betekenis\%20sector\%20EV\%20Nederland\%202020-2022.pdf}
\source{nlseba}{Netherlands Enterprise Agency}{Subsidieregeling Emissieloze Bedrijfsautos (SEBA)}{https://www.rvo.nl/subsidies-financiering/seba}
\end{thebibliography}
\end{document}
